\documentclass[11pt,letterpaper]{article}
\usepackage[margin=1.2in]{geometry}
\usepackage[scaled=0.95]{helvet}
\renewcommand{\familydefault}{\sfdefault}
\usepackage[T1]{fontenc}
\usepackage{microtype}
\usepackage{xcolor}
\definecolor{softblack}{HTML}{1A1A1A}
\definecolor{mutedgray}{HTML}{555555}
\AtBeginDocument{\color{softblack}}
\setlength{\parskip}{6pt}
\setlength{\parindent}{0pt}
\pagestyle{empty}

\begin{document}
\noindent\textbf{Cover letter --- \emph{American Economic Journal: Macroeconomics}}\\[0.4em]
\textcolor{mutedgray}{\rule{0.22\linewidth}{0.4pt}}

\bigskip
\noindent Editor\\
\textit{American Economic Journal: Macroeconomics}

\bigskip

\noindent Dear Editor,

\bigskip

\noindent We are pleased to submit ``Volcanic Forcing and the Pre-Industrial Malthusian Trap: Pathway-Heterogeneous Demographic Response and the Long Shadow of Seasonality'' for consideration at the \emph{American Economic Journal: Macroeconomics}.

The paper estimates how pre-industrial economies absorbed climate shocks along demographic and agricultural-land margins, using volcanic eruptions as exogenous forcing and a new monthly paleo-economic panel that we assemble from the ModE-RA paleo-reanalysis, HYDE 3.5 land-use grids, the eVolv2k ice-core volcanic sulfur record, the Brecke conflict catalogue, and the Allen-Nuffield wage archive (196 countries, 3{,}187 sub-national units, 1421--2008~CE). The headline finding is a sharply pathway-heterogeneous Malthusian demographic margin: crop-dominant late and pastoral/mixed late countries lose population at $-9.7$ and $-5.5\!\cdot\!10^{-5}$ per Tg of stratospheric sulfur injection respectively, with the magnitudes implying $\sim 1\%$ less annualised population growth in a typical Tambora-decade exposure and $\sim 10\%$ cumulative loss for a country exposed to a Tambora-class shock without intervening recovery. Both coefficients are Bonferroni-significant across twelve pathway-by-equation tests, survive re-clustering pathways on climate primitives orthogonal to HYDE outcomes, survive a Webb wild-cluster bootstrap, and survive restriction to the post-1700 panel where HYDE's back-projection is most reliable.

The macroeconomic substance the paper delivers, in our reading, sits in three findings of direct relevance to this journal.

\emph{First}, the Malthusian channel has a quantifiable welfare arithmetic. Combining the demographic-margin coefficient with the Allen real-wage Malthus coefficient ($\hat\beta_{\ln P} = -0.325$ on the city-buffer specification, 1421--1850), a 30~Tg Tambora-class decade implies a 3\% cumulative population decline in crop-dominant late countries and a corresponding $+0.95\%$ real-wage rise for the survivors---the canonical Malthusian welfare paradox at climate-shock magnitudes. A calibrated three-equation structural model traces the cumulative impact of the 1500--1890 VSSI trajectory: in the no-volcanism counterfactual, crop-dominant late population grows by an additional $+212\%$ over the four centuries (with the surviving per-capita welfare loss of $-0.32$ log-units, the Malthusian-paradox demographic dividend the literature predicts).

\emph{Second}, the pre-industrial climate dimension that selected which agricultural system a society adopted leaves measurable traces in modern demographic trajectories: pre-industrial inter-annual climate volatility 1421--1750 is the single strongest cross-country predictor of modern (1955--1960 to 2015--2020) population growth at $R^2 = 0.45$, surviving an extended deep-determinants battery (Neolithic distance, migratory distance, ruggedness, land area, log GDP per capita) and an absolute-latitude control, and failing a placebo when the same volatility measure is computed over the modern 1950--2008 reanalysis period. The macroeconomic implication is that pre-industrial climate regimes shape modern demographic ceilings through institutional and technological channels we cannot disentangle in a cross-country panel, but the geographic component of the channel is identified within-country.

\emph{Third}, we test the Boserupian intensification hypothesis at cross-country pre-industrial scale and report a measurement-driven null on HYDE land-use data plus a small but statistically robust signal on the independent KK10 reconstruction. The HYDE cropland-share specification is artifactually positive (driven by colonial-transition encoding in small-island countries); the cropland-area specification, while statistically robust, is fully mediated by the demographic margin (HYDE uses population as an input). On KK10, which is population-independent by construction, the high-density intensive pathway shows a Boserupian coefficient of $+4.3\!\cdot\!10^{-5}$ per Tg ($p=0.047$, surviving the demographic-margin control)---the pathway where the demographic response is weakest also shows the strongest intensification response, qualitatively consistent with channel substitution. The macroeconomic implication is that the Boserupian channel exists at pre-industrial scale, but its magnitude is two orders below what the HYDE-cropland-share specifications in the existing climate-economics literature had suggested.

The methodological documentation of HYDE's two-axis identification problem (share-specification denominator instability + area-specification demographic mediation) is itself a contribution to the empirical climate-economics literature this journal has shaped. The KK10 cross-validation strategy generalises directly.

We see the paper as a contribution to three lines of work \emph{AEJ:Macro} has been central to: (i) the macroeconomics of climate shocks at pre-industrial as well as modern scale (Burke, Hsiang and Miguel 2015; Dell, Jones and Olken 2012); (ii) the long-run consequences of agricultural pathways and climate endowments for modern development (Ashraf and Galor 2011; Nunn and Qian 2011); and (iii) the structural dynamics of the Malthusian regime and its eventual escape, including the demographic transition signature we recover in the within-country regressions.

The paper is coauthored by two researchers at the Centro de Investigaci\'on Econ\'omica, ITAM, and Universidade de Vigo / RGEA, Spain. We declare no conflicts of interest. The manuscript has not been previously published and is not under consideration elsewhere. A replication package with the ModE-RA aggregation code, the KK10 country aggregation script, the country and sub-national parquet panels, and all analysis scripts will be deposited on OSF upon acceptance.

\bigskip

\noindent Sincerely,\\[0.4em]
Jorge Alonso Ortiz \\
\small\textcolor{mutedgray}{Centro de Investigaci\'on Econ\'omica, ITAM\quad\textbar\quad jorge.alonso@itam.mx}

\bigskip

\noindent Jos\'e-Mar\'ia Da-Rocha \\
\small\textcolor{mutedgray}{Universidade de Vigo \& RGEA, Spain\quad\textbar\quad darocha@uvigo.es}

\bigskip

\textcolor{mutedgray}{\rule{0.22\linewidth}{0.4pt}}\\[0.2em]
{\footnotesize\textcolor{mutedgray}{Suggested referees (independent of the authors):}\\
M.\ Burke (Stanford, climate-economy);
T.\ B.\ Andersen (Copenhagen, climate-growth);
B.\ Olken (MIT, climate-development);
A.\ Matranga (NES, seasonality-agriculture);
M.\ Lewis (Brown, paleo-economy).}

\end{document}
